\section{Introduction}
\label{sec:intro}

The Vesuvius Challenge asks for text read from X-ray micro-CT volumes of carbonised scrolls without opening them \cite{seales2023,scrollprize}. The published route has three stages: segment a surface mesh that follows one winding of the rolled sheet, render the CT intensity in a thin slab around that mesh into a flat image stack, and run an ink-detection network on the stack. Stage one is the bottleneck. A tracer grown from a seed point will, with some probability per unit area, cross the interstitial gap and continue on the neighbouring winding. The resulting mesh is smooth, its rendering looks like papyrus, and only the fibre pattern or a cross-section reveals that the surface now covers two different sheets. The organisers' own criterion for unwrapping work is that \emph{visible papyrus fibres do not jump across sheets} \cite{scrollprize}.

This note has two aims. The first is documentary: a complete, headless, timed run of the open-source tools on PHerc1447 from a laptop with no CUDA device, including the failures, so that the next person starts from a working baseline. The second is methodological. Both quality checks introduced here come from physics rather than from image processing, and both are stated with their derivations so that they can be criticised at the level of the equations.

\paragraph{Developability.} A sheet of papyrus is, to the precision of the scan, an isometric embedding of a flat rectangle. Its Gaussian curvature vanishes, so parallel transport of a tangent vector around any closed loop on it returns the vector unchanged. This is the same statement as the Thomas precession of a gyroscope or the Foucault pendulum: the holonomy angle equals the integrated curvature enclosed by the loop. \Cref{sec:holonomy} turns this into a per-loop test on a quad mesh via the discrete Gauss--Bonnet theorem, and reports a negative result with a reason: a sheet switch is a change of \emph{which} surface the mesh follows, not a change of the mesh's intrinsic geometry.

\paragraph{Tetratic order.} The papyrus weave consists of two orthogonal families of fibres, recto and verso. An orientation field with this symmetry is a tetratic, the four-fold analogue of a nematic liquid crystal \cite{degennes}. \Cref{sec:tetratic} defines the order parameter $\psifour$ on tiles of a flattened render and the disclination charge of its phase, and shows that $|\psifour|$ drops from \numrange{0.6}{0.8} on clean weave to \numrange{0.1}{0.3} across a sheet switch, while the structure-tensor coherence that one would normally use stays flat.

\paragraph{Ink detection on a laptop.} \Cref{sec:ink} describes inference with the released \code{ink\_9um} and Hecate checkpoints on the Apple Metal backend, a positive control on PHerc.172 segment w087 against the team's published ink map, a scan of the model's 17-layer depth window through a 28-layer render, and a failed lock-in idea. All of it is scored with Pearson correlation and a Mann--Whitney AUC, with block-bootstrap intervals, so that the differences between settings can be read as evidence rather than as impressions.

The document is a working draft. Sections are written so that results can be appended as the cluster runs described in \cref{sec:discussion} complete.
